\chapter{Speicherung und Metadaten}
\label{ch:map/storage}

In diesem Kapitel lernst du, wie du Daten speichern kannst.

\section{Metadaten}\label{metadaten}

\subsection{Was sind Metadaten?}\label{was-sind-metadaten}

Metadaten\index{metadaten@Metadaten} sind in Luanti ein
Schlüssel-Wert-Speicher, mit dem du eigene Daten an etwas anhängst. Du
kannst damit Informationen zu einem Node, einem Spieler oder einem
ItemStack\index{itemstack@\texttt{ItemStack}} speichern.

Alle Arten von Metadaten verwenden genau dieselbe API. Metadaten
speichern Werte als Strings, es gibt aber eine Reihe von Methoden, um
andere einfache Datentypen umzuwandeln und zu speichern.

Manche Schlüssel haben eine besondere Bedeutung. So speichert
\texttt{infotext} in Node-Metadaten den Tooltip, der erscheint, wenn man
mit dem Fadenkreuz auf den Node zeigt. Um Konflikte mit anderen Mods zu
vermeiden, solltest du für Schlüssel die übliche Namensraum-Konvention
verwenden: \texttt{modname:keyname}. Ausnahme sind gebräuchliche Angaben
wie der Name des Besitzers, der als \texttt{owner} gespeichert wird.

Metadaten sind Daten über Daten. Die Daten selbst, etwa der Typ eines
Nodes oder die Anzahl in einem Stack, sind keine Metadaten.

\subsection{Ein Metadaten-Objekt
holen}\label{ein-metadaten-objekt-holen}

Kennst du die Position eines Nodes, kannst du seine Metadaten abrufen:

\begin{Shaded}
\begin{Highlighting}[]
\KeywordTok{local} \VariableTok{meta} \OperatorTok{=} \VariableTok{core}\OperatorTok{.}\NormalTok{get\_meta}\OperatorTok{(\{} \VariableTok{x} \OperatorTok{=} \DecValTok{1}\OperatorTok{,} \VariableTok{y} \OperatorTok{=} \DecValTok{2}\OperatorTok{,} \VariableTok{z} \OperatorTok{=} \DecValTok{3} \OperatorTok{\})}
\end{Highlighting}
\end{Shaded}

Spieler- und ItemStack-Metadaten erhältst du mit
\texttt{get\_\allowbreak{}meta()}:

\begin{Shaded}
\begin{Highlighting}[]
\KeywordTok{local} \VariableTok{pmeta} \OperatorTok{=} \VariableTok{player}\OperatorTok{:}\NormalTok{get\_meta}\OperatorTok{()}
\KeywordTok{local} \VariableTok{imeta} \OperatorTok{=} \VariableTok{stack}\OperatorTok{:}\NormalTok{get\_meta}\OperatorTok{()}
\end{Highlighting}
\end{Shaded}

\subsection{Lesen und Schreiben}\label{lesen-und-schreiben}

Meist liest und schreibst du Metadaten mit den Methoden
\texttt{get\_\allowbreak{}<typ>()} und
\texttt{set\_\allowbreak{}<typ>()}. Da Metadaten Strings speichern,
setzen und lesen die String-Methoden den Wert direkt.

\begin{Shaded}
\begin{Highlighting}[]
\FunctionTok{print}\OperatorTok{(}\VariableTok{meta}\OperatorTok{:}\NormalTok{get\_string}\OperatorTok{(}\StringTok{"foo"}\OperatorTok{))} \CommentTok{{-}{-}\textgreater{} ""}
\VariableTok{meta}\OperatorTok{:}\NormalTok{set\_string}\OperatorTok{(}\StringTok{"foo"}\OperatorTok{,} \StringTok{"bar"}\OperatorTok{)}
\FunctionTok{print}\OperatorTok{(}\VariableTok{meta}\OperatorTok{:}\NormalTok{get\_string}\OperatorTok{(}\StringTok{"foo"}\OperatorTok{))} \CommentTok{{-}{-}\textgreater{} "bar"}
\end{Highlighting}
\end{Shaded}

Alle typisierten Getter geben einen neutralen Standardwert zurück, wenn
der Schlüssel nicht existiert, etwa
\texttt{\textquotedbl{}\textquotedbl{}} oder \texttt{0}. Mit
\texttt{get()} erhältst du einen String oder \texttt{nil}.

Metadaten-Objekte sind Referenzen; Änderungen landen also automatisch an
der Quelle. ItemStacks sind dagegen keine Referenzen; nach einer
Änderung musst du den ItemStack wieder ins
Inventar\index{inventar@Inventar} setzen.

Die Getter und Setter für andere Typen wandeln von und nach String um:

\begin{Shaded}
\begin{Highlighting}[]
\FunctionTok{print}\OperatorTok{(}\VariableTok{meta}\OperatorTok{:}\NormalTok{get\_int}\OperatorTok{(}\StringTok{"count"}\OperatorTok{))} \CommentTok{{-}{-}\textgreater{} 0}
\VariableTok{meta}\OperatorTok{:}\NormalTok{set\_int}\OperatorTok{(}\StringTok{"count"}\OperatorTok{,} \DecValTok{3}\OperatorTok{)}
\FunctionTok{print}\OperatorTok{(}\VariableTok{meta}\OperatorTok{:}\NormalTok{get\_int}\OperatorTok{(}\StringTok{"count"}\OperatorTok{))} \CommentTok{{-}{-}\textgreater{} 3}
\FunctionTok{print}\OperatorTok{(}\VariableTok{meta}\OperatorTok{:}\NormalTok{get\_string}\OperatorTok{(}\StringTok{"count"}\OperatorTok{))} \CommentTok{{-}{-}\textgreater{} "3"}
\end{Highlighting}
\end{Shaded}

\subsection{Besondere Schlüssel}\label{besondere-schluxfcssel}

\texttt{infotext} zeigt in Node-Metadaten einen Tooltip an, wenn das
Fadenkreuz auf den Node zeigt. Das ist nützlich, um Besitzer oder
Zustand eines Nodes anzuzeigen.

\texttt{description} überschreibt in ItemStack-Metadaten die
Beschreibung, die beim Überfahren des Stacks im Inventar erscheint.
Farben kannst du mit
\texttt{core.\allowbreak{}colorize()}\index{core@\texttt{core}!colorize@\texttt{core.\allowbreak{}colorize}}
einfügen.

\texttt{owner} ist ein gebräuchlicher Schlüssel für den Namen des
Spielers, dem das Item oder der Node gehört.

\subsection{Tabellen speichern}\label{tabellen-speichern}

Tabellen müssen in Strings umgewandelt werden, bevor du sie speichern
kannst. Luanti bietet dafür zwei Formate: Lua und JSON.

Das Lua-Format ist meist deutlich schneller und entspricht der
Tabellenschreibweise von Lua. JSON ist ein verbreiteterer Standard,
besser strukturiert und gut geeignet, wenn du Daten mit einem anderen
Programm austauschen musst. Für JSON gibt es
\texttt{core.\allowbreak{}write\_\allowbreak{}json()}\index{core@\texttt{core}!write_json@\texttt{core.\allowbreak{}write\_\allowbreak{}json}}
und
\texttt{core.\allowbreak{}parse\_\allowbreak{}json()}\index{core@\texttt{core}!parse_json@\texttt{core.\allowbreak{}parse\_\allowbreak{}json}}.

\begin{Shaded}
\begin{Highlighting}[]
\KeywordTok{local} \VariableTok{data} \OperatorTok{=} \OperatorTok{\{} \VariableTok{username} \OperatorTok{=} \StringTok{"player1"}\OperatorTok{,} \VariableTok{score} \OperatorTok{=} \DecValTok{1234} \OperatorTok{\}}
\VariableTok{meta}\OperatorTok{:}\NormalTok{set\_string}\OperatorTok{(}\StringTok{"foo"}\OperatorTok{,} \VariableTok{core}\OperatorTok{.}\NormalTok{serialize}\OperatorTok{(}\VariableTok{data}\OperatorTok{))}

\VariableTok{data} \OperatorTok{=} \VariableTok{core}\OperatorTok{.}\NormalTok{deserialize}\OperatorTok{(}\VariableTok{meta}\OperatorTok{:}\NormalTok{get\_string}\OperatorTok{(}\StringTok{"foo"}\OperatorTok{))}
\end{Highlighting}
\end{Shaded}

\subsection{Private Metadaten}\label{private-metadaten}

Standardmäßig werden alle Node-Metadaten an den Client geschickt. Um das
zu verhindern, kannst du Schlüssel als privat markieren:

\begin{Shaded}
\begin{Highlighting}[]
\VariableTok{meta}\OperatorTok{:}\NormalTok{set\_string}\OperatorTok{(}\StringTok{"secret"}\OperatorTok{,} \StringTok{"asd34dn"}\OperatorTok{)}
\VariableTok{meta}\OperatorTok{:}\NormalTok{mark\_as\_private}\OperatorTok{(}\StringTok{"secret"}\OperatorTok{)}
\end{Highlighting}
\end{Shaded}

\subsection{Lua-Tabellen}\label{lua-tabellen}

Mit \texttt{to\_\allowbreak{}table} und
\texttt{from\_\allowbreak{}table} wandelst du Metadaten in Lua-Tabellen
um und zurück:

\begin{Shaded}
\begin{Highlighting}[]
\KeywordTok{local} \VariableTok{tmp} \OperatorTok{=} \VariableTok{meta}\OperatorTok{:}\NormalTok{to\_table}\OperatorTok{()}
\VariableTok{tmp}\OperatorTok{.}\VariableTok{foo} \OperatorTok{=} \StringTok{"bar"}
\VariableTok{meta}\OperatorTok{:}\NormalTok{from\_table}\OperatorTok{(}\VariableTok{tmp}\OperatorTok{)}
\end{Highlighting}
\end{Shaded}

\section{Mod-Speicher}\label{mod-speicher}

Der Mod-Speicher\index{mod-speicher@Mod-Speicher} (\emph{mod storage})
verwendet genau dieselbe API wie Metadaten, auch wenn er technisch
gesehen keine Metadaten sind. Er gehört jeweils zu einem Mod und kann
nur während der Ladezeit abgerufen werden, damit Luanti weiß, welcher
Mod ihn anfordert.

\begin{Shaded}
\begin{Highlighting}[]
\KeywordTok{local} \VariableTok{storage} \OperatorTok{=} \VariableTok{core}\OperatorTok{.}\NormalTok{get\_mod\_storage}\OperatorTok{()}
\end{Highlighting}
\end{Shaded}

Jetzt kannst du den Speicher genau wie Metadaten verwenden:

\begin{Shaded}
\begin{Highlighting}[]
\VariableTok{storage}\OperatorTok{:}\NormalTok{set\_string}\OperatorTok{(}\StringTok{"foo"}\OperatorTok{,} \StringTok{"bar"}\OperatorTok{)}
\end{Highlighting}
\end{Shaded}

\section{Datenbanken}\label{datenbanken}

Wird ein Mod wahrscheinlich auf Servern eingesetzt und speichert er
viele Daten, lohnt es sich, eine Datenbank als Speicher anzubieten. Mach
das optional, indem du trennst, \emph{wie} die Daten gespeichert werden
und \emph{wo} sie verwendet werden.

\begin{Shaded}
\begin{Highlighting}[]
\KeywordTok{local} \VariableTok{backend}
\ControlFlowTok{if} \VariableTok{use\_database} \ControlFlowTok{then}
    \VariableTok{backend} \OperatorTok{=} \FunctionTok{dofile}\OperatorTok{(}
        \VariableTok{core}\OperatorTok{.}\NormalTok{get\_modpath}\OperatorTok{(}\StringTok{"mymod"}\OperatorTok{)} \OperatorTok{..} \StringTok{"/backend\_sqlite.lua"}
    \OperatorTok{)}
\ControlFlowTok{else}
    \VariableTok{backend} \OperatorTok{=} \FunctionTok{dofile}\OperatorTok{(}
        \VariableTok{core}\OperatorTok{.}\NormalTok{get\_modpath}\OperatorTok{(}\StringTok{"mymod"}\OperatorTok{)} \OperatorTok{..} \StringTok{"/backend\_storage.lua"}
    \OperatorTok{)}
\ControlFlowTok{end}

\VariableTok{backend}\OperatorTok{.}\NormalTok{get\_foo}\OperatorTok{(}\StringTok{"a"}\OperatorTok{)}
\VariableTok{backend}\OperatorTok{.}\NormalTok{set\_foo}\OperatorTok{(}\StringTok{"a"}\OperatorTok{,} \OperatorTok{\{} \VariableTok{score} \OperatorTok{=} \DecValTok{3} \OperatorTok{\})}
\end{Highlighting}
\end{Shaded}

Die Datei \texttt{backend\_\allowbreak{}storage.\allowbreak{}lua}
enthält dann eine Umsetzung mit dem Mod-Speicher:

\begin{Shaded}
\begin{Highlighting}[]
\KeywordTok{local} \VariableTok{storage} \OperatorTok{=} \VariableTok{core}\OperatorTok{.}\NormalTok{get\_mod\_storage}\OperatorTok{()}
\KeywordTok{local} \VariableTok{backend} \OperatorTok{=} \OperatorTok{\{\}}

\KeywordTok{function} \VariableTok{backend}\OperatorTok{.}\NormalTok{set\_foo}\OperatorTok{(}\VariableTok{key}\OperatorTok{,} \VariableTok{value}\OperatorTok{)}
    \VariableTok{storage}\OperatorTok{:}\NormalTok{set\_string}\OperatorTok{(}\VariableTok{key}\OperatorTok{,} \VariableTok{core}\OperatorTok{.}\NormalTok{serialize}\OperatorTok{(}\VariableTok{value}\OperatorTok{))}
\KeywordTok{end}

\KeywordTok{function} \VariableTok{backend}\OperatorTok{.}\NormalTok{get\_foo}\OperatorTok{(}\VariableTok{key}\OperatorTok{)}
    \ControlFlowTok{return} \VariableTok{core}\OperatorTok{.}\NormalTok{deserialize}\OperatorTok{(}\VariableTok{storage}\OperatorTok{:}\NormalTok{get\_string}\OperatorTok{(}\VariableTok{key}\OperatorTok{))}
\KeywordTok{end}

\ControlFlowTok{return} \VariableTok{backend}
\end{Highlighting}
\end{Shaded}

\texttt{backend\_\allowbreak{}sqlite.\allowbreak{}lua} macht dasselbe,
verwendet aber statt des Mod-Speichers die Lua-Bibliothek
\emph{lsqlite3}.

Für eine Datenbank wie SQLite brauchst du eine unsichere Umgebung
(\emph{insecure environment}). Das ist eine Tabelle, die nur Mods zur
Verfügung steht, die der Benutzer ausdrücklich freigegeben hat. Sie
enthält eine weniger eingeschränkte Kopie der Lua-API, die von
bösartigen Mods missbraucht werden könnte. Mehr dazu im Kapitel
Sicherheit~(Kapitel~\ref{ch:quality/security}).

\begin{Shaded}
\begin{Highlighting}[]
\KeywordTok{local} \VariableTok{ie} \OperatorTok{=} \VariableTok{core}\OperatorTok{.}\NormalTok{request\_insecure\_environment}\OperatorTok{()}
\ControlFlowTok{if} \KeywordTok{not} \VariableTok{ie} \ControlFlowTok{then}
    \FunctionTok{error}\OperatorTok{(}\StringTok{"Bitte mymod zu secure.trusted\_mods hinzufügen"}\OperatorTok{)}
\ControlFlowTok{end}

\KeywordTok{local} \VariableTok{\_sql} \OperatorTok{=} \VariableTok{ie}\OperatorTok{.}\NormalTok{require}\OperatorTok{(}\StringTok{"lsqlite3"}\OperatorTok{)}
\CommentTok{{-}{-} globale sqlite3{-}Bibliothek vor anderen Mods verbergen}
\ControlFlowTok{if} \VariableTok{sqlite3} \ControlFlowTok{then}
    \VariableTok{sqlite3} \OperatorTok{=} \KeywordTok{nil}
\ControlFlowTok{end}
\end{Highlighting}
\end{Shaded}

SQL oder die Bibliothek lsqlite3 zu erklären, würde den Rahmen dieses
Buches sprengen.

\section{Was soll ich verwenden?}\label{was-soll-ich-verwenden}

Welche Methode passt, hängt davon ab, worauf sich die Daten beziehen,
wie sie aufgebaut sind und wie groß sie sind. Als Faustregel: Kleine
Datenmengen gehen bis 10 KB, mittlere bis 10 MB, alles darüber ist groß.

Node-Metadaten sind eine gute Wahl für Daten, die zu einem Node gehören.
Mittlere Datenmengen lassen sich recht effizient speichern, wenn du sie
privat markierst.

Item-Metadaten solltest du nur für kleine Datenmengen verwenden, denn
sie werden immer an den Client geschickt. Außerdem werden die Daten
jedes Mal kopiert, wenn der Stack bewegt oder aus Lua darauf zugegriffen
wird.

Der Mod-Speicher eignet sich für mittlere Datenmengen; große Datenmengen
zu schreiben, kann aber ineffizient sein. Für große Datenmengen nimmst
du besser eine Datenbank, damit nicht bei jedem Speichern alle Daten
geschrieben werden müssen.

Datenbanken kommen nur für Server in Frage, weil der Mod für die
unsichere Umgebung freigegeben werden muss. Für große Datenmengen sind
sie gut geeignet.

\section{Jetzt bist du dran}\label{jetzt-bist-du-dran}

\begin{itemize}
\tightlist
\item
  Erstelle einen Node, der verschwindet, nachdem er fünfmal geschlagen
  wurde. (Verwende
  \texttt{on\_\allowbreak{}punch}\index{callback@Callback!on_punch@\texttt{on\_\allowbreak{}punch}}
  in der Node-Definition und
  \texttt{core.\allowbreak{}set\_\allowbreak{}node}\index{core@\texttt{core}!set_node@\texttt{core.\allowbreak{}set\_\allowbreak{}node}}.)
\end{itemize}
